\documentclass{article}
\usepackage[UTF8]{ctex}
\usepackage{amsmath,amssymb,amsfonts}
\usepackage{enumitem}
\usepackage{geometry}
\geometry{a4paper, margin=1in}
% === 错题本排版优化（不参与题目标识） ===
\setlength{\parindent}{0pt}
\setlength{\parskip}{0.35em plus 0.10em minus 0.08em}
\setlength{\emergencystretch}{3em}
\setlist{itemsep=0.25em, topsep=0.40em, parsep=0pt}
\allowdisplaybreaks[2]
% === 错题本排版优化结束 ===
\begin{document}
\section*{华中师范大学2025年数学分析真题}
\section*{139．华中师范大学2025年数学分析试卷}

\subsection{题目1}
1．求极限 \(\lim _{R \rightarrow+\infty} \iint_{D}\left(x^{2}+y^{2}\right) e^{-\left(x^{2}+y^{2}\right)} \mathrm{d} x \mathrm{~d} y\) ，其中区域 \(D=\{(x, y) \| x|\leq R,|y| \leq R\}\) ．

\subsection{题目2}
2．设 \(\varphi(t)=\int_{0}^{1} \ln \sqrt{x^{2}+t^{2}} \mathrm{~d} x(0 \leq t \leq 1)\) ，求 \(\varphi\) 在 \(t=0\) 处的单侧导数。

\subsection{题目3}
3．计算 \(J=\int_{0}^{\pi} \frac{x \sin x}{1+\cos ^{2} x} \mathrm{~d} x\) ．

\subsection{题目4}
4 ．求全微分 \((y+z) \mathrm{d} x+(z+x) \mathrm{d} y+(x+y) \mathrm{d} z\) 的原函数 \(u(x, y, z)\) ．

\subsection{题目5}
5．计算 \(\iint_{S} z^{3} \mathrm{~d} x \mathrm{~d} y\) ，其中 \(S\) 是椭球面 \(\frac{x^{2}}{a^{2}}+\frac{y^{2}}{b^{2}}+\frac{z^{2}}{c^{2}}=1(a, b, c>0)\) 的上半部分，并取外侧．

二．（10分）用数列极限的柯西准则证明 \(\lim _{n \rightarrow \infty} \sin n\) 不存在．\\

\subsection{题目3}
原题拆分片段：
三（15 分）设不含原点的开区域 \(D \subset \mathbb{R}^{2}\) 上的二元可微函数 \(F\) 在直角坐标系下可表示为

\[
z=F(x, y)=f(x) g(y), f(x) \neq 0, g(y) \neq 0 .
\]

它在极坐标系可表示为

\[
z=F(r \cos \theta, r \sin \theta)=h(r) .
\]

求二元函数 \(z=F(x, y)\) 的显式表达式．\\

\subsection{题目6}
13．（15 分）考察级数 \(\sum_{n=1}^{\infty}\left(1+\frac{1}{2}+\cdots+\frac{1}{n}\right) \frac{\sin n}{n}\) 的敛散性。

五．（15 分）设函数 \(f\) 在 \(\mathbb{R}\) 上的二阶导数 \(f^{\prime \prime}(x)>0\) ，又设存在点 \(x_{0} \in \mathbb{R}\) ，使得 \(f\left(x_{0}\right)<0\) ，且

\[
\lim _{x \rightarrow-\infty} f^{\prime}(x)=\alpha<0, \lim _{x \rightarrow+\infty} f^{\prime}(x)=\beta>0 .
\]

证明：\(f(x)\) 有且仅有 2 个零点．

六．（10 分）设函数 \(f\) 在有限区间 \([a, b]\) 上连续，且对任意的 \([\alpha, \beta] \subset[a, b]\) ，有

\[
\left|\int_{\alpha}^{\beta} f(x) \mathrm{d} x\right| \leq M|\beta-\alpha|^{s+1}
\]

其中 \(M, s\) 为正常数．证明：\(f(x) \equiv 0, x \in[a, b]\) ．\\

\subsection{题目9}
原题边界“L”：（15 分）设函数 \(f\) 在任意有界区间上可积，并且 \(\lim _{x \rightarrow+\infty} f(x)=\alpha\)（ \(\alpha\) 为有限数）．证明：

\[
\lim _{y \rightarrow 0^{+}} y \int_{0}^{+\infty} e^{-y x} f(x) \mathrm{d} x=\alpha
\]

\subsection{题目9}
原题边界“A”：（20 分）给定级数 \(\sum_{n=1}^{\infty} \frac{a_{n}}{n^{x}}\) ，其中 \(a_{n}\) 均为实数。证明：\\
（1）（8分）若 \(\sum_{n=1}^{\infty} \frac{a_{n}}{n^{\lambda}}\) 收敛，则对任意的 \(x>\lambda, \sum_{n=1}^{\infty} \frac{a_{n}}{n^{x}}\) 收敛。\\
（2）（12 分）存在 \(r \in[-\infty,+\infty]\) ，使得 \(x>r\) 时，\(\sum_{n=1}^{\infty} \frac{a_{n}}{n^{x}}\) 收敛；\(x<r\) 时，\(\sum_{n=1}^{\infty} \frac{a_{n}}{n^{x}}\) 发散．
\end{document}
